%% kapa: exam=1 round=32 year=2021 period=1 form=A subject=economics subject_ko=경제학원론 seq=13 number=53
%% source: Q-Net 감정평가사 기출문제 (cst00309) artlSeq=5208397 file=extracted/1008/5208397/A형/2021년도 제32회 감정평가사 1차 1교시 A형.pdf
%% conversion: kordoc 4.24.51 -> tools/gen1.py (auto); review=proofread-sonnet; flags=none
\begin{question}{53}
기업 甲과 乙만 있는 상품시장에서 두 기업이 꾸르노(Cournot) 모형에 따라 행동하는 경우에 관한 설명으로 옳은 것을 모두 고른 것은? (단, 생산기술은 동일하다.)
\begin{framedtext}
ㄱ. 甲은 乙이 생산량을 결정하면 그대로 유지될 것이라고 추측한다. \par
ㄴ. 甲과 乙은 생산량 결정에서 서로 협력한다. \par
ㄷ. 甲, 乙 두 기업이 완전한 담합을 이루는 경우와 꾸르노 균형의 결과는 동일하다. \par
ㄹ. 추가로 기업이 시장에 진입하는 경우 균형가격은 한계비용에 접근한다. \par
\end{framedtext}
\begin{choices}
  \item ㄱ, ㄷ
  \item ㄱ, ㄹ
  \item ㄴ, ㄷ
  \item ㄴ, ㄹ
  \item ㄷ, ㄹ
\end{choices}
\end{question}
\answer{2}
